\chapter*{ Введение }
\label{ch:ch0}

Изучение моделей в статистической физике позволяет  найти связь микроскопических свойств системы с термодинамическими характеристиками. 
 Статистическая физика изучает коллективное поведение больших систем через упрощенные эффективные модели, так как полное микроскопическое описание с учётом всех степеней свободы вычислительно сложно и избыточно по информации.
 
 Одним из основных интересов изучения сложных систем является поиск критических точек, в которых, как правило реализуется фазовый переход --- смена состояние системы \cite{cardy1996scaling}.
Вблизи фазовых переходов макроскопические свойства системы  определяются небольшим числом признаков (например, симметрией, размерностью пространства и дальностью взаимодействий), что приводит к формированию классов универсальности. 
В окрестности критических точек микродетали сложных систем становятся нерелевантными, поэтому разные физические системы демонстрируют одни и те же критические характеристики. Классический пример --- модель Изинга, которая относится к тому же классу универсальности, что и решеточный газ (который является приближением к реальной жидкости) \cite{PELISSETTO2002549}. 


Фазовые переходы бывают различной природы. Порядок перехода определяет, какие именно универсальные законы применимы и какие анзацы корректны для экстраполяции к термодинамическому пределу.


Переходы первого рода характеризуются разрывом первой производной свободной энергии и сосуществованием фаз. 

Фазовые переходы второго рода (или непрерывные) характеризуются отсутствием  разрыва параметра порядка, однако корреляционная длина и флуктуации расходятся вблизи критической точки. 
Для таких переходов определяются критические индексы и выполняются скейлинговые соотношения, а их значения зависят только от универсальных признаков системы (симметрии, размерности и дальности взаимодействий), но не от микроскопических деталей. Примерами систем, реализующих переход второго рода являются модели Изинга на квадратной и кубической решетках, а также модель XY на кубической решетке \cite{PELISSETTO2002549}. 

Особый тип непрерывных переходов реализуется в двумерных системах с непрерывной симметрией, например в XY-модели  \cite{kosterlitz1973ordering}. Это переход Березинского–Костерлица–Таулеса (БКТ), в котором нет локального параметра порядка: фазовый переход связан с развязыванием пар вихрь–антивихрь (vortices ), а корреляционная функция меняет характер от степенного убывания к экспоненциальному. 
Для БКТ-перехода характерны экспоненциальное расходимость корреляционной длины. 




%В большинстве интересных моделей точные решения отсутствуют, а асимптотические методы дают лишь ограниченные предсказания.